\chapter{Results and Discussion}
\label{ch:results}

\section{C++ Reference Quantization Sweep}

The C++ reference sweep evaluates several quantization configurations. \Cref{tab:cpp-sweep} summarizes the observed behavior.

\begin{table}[h]
\centering
\caption{Summary of C++ reference quantization sweep.}
\label{tab:cpp-sweep}
\begin{tabular}{llll}
\toprule
Configuration & Main override & Observation & Interpretation \\
\midrule
\texttt{fake} & none & Nearly identical to float & Reference upper bound \\
\texttt{all\_int16} & \texttt{ssm\_state=int16, gate=int16} & Similar task error to fake & int16 is safe \\
\texttt{mixed\_candidate} & \texttt{dt\_proj=int8, ssm\_state=int8, gate=int8} & Similar to int16 & Good mixed candidate \\
\texttt{patch\_i8\_on\_i16} & \texttt{patch\_embedding=int8} & Large degradation & Patch embedding is int8-sensitive \\
\texttt{inproj\_i8\_on\_i16} & \texttt{in\_proj=int8} & Large degradation & Input projection is int8-sensitive \\
\texttt{outproj\_i8\_on\_i16} & \texttt{out\_proj=int8} & Moderate degradation & Output projection needs care \\
\bottomrule
\end{tabular}
\end{table}

The sweep indicates that the overall int16 reference preserves the float model behavior, while selected projection layers are sensitive to int8 quantization.

\section{Stage-Level Error Localization}

Stage-level comparison uses:

\begin{equation}
  \mathrm{MAE}_{\mathrm{stage}} =
  \mathrm{mean}\left(|x_{\mathrm{cpp\_like}} - x_{\mathrm{hw\_like}}|\right).
\end{equation}

This comparison locates the first stage where RTL-like fixed-point semantics diverge from the C++ reference semantics.

\section{Selective Scan Error in Case 04}

In case 04, the fixed-Q8.8 state representation produced a high selective scan MAE:

\begin{equation}
  \mathrm{MAE}_{\mathrm{selective\_scan, fixed}} = 0.121098.
\end{equation}

Earlier stages such as normalization, input projection, and sigmoid had smaller errors. Therefore, the selective scan state representation became the main suspect.

\section{Scaled-State Quantization}

The scaled-state method uses per-state-channel scale values generated by Python and consumed by the hardware-like and RTL paths. The selective scan state is stored as a scaled int16 value and converted back to Q8.8 for downstream stages.

After applying scaled-state quantization, the average selective scan MAE decreased to:

\begin{equation}
  \mathrm{MAE}_{\mathrm{selective\_scan, scaled}} = 0.090323.
\end{equation}

The improvement is:

\begin{equation}
  \frac{0.121098 - 0.090323}{0.121098} \approx 25.41\%.
\end{equation}

\section{Error Propagation}

The selective scan output feeds the element-wise gate and then the output projection. Therefore, improvements in the state representation propagate to later stages, but the final block output improvement is smaller because residual addition and projection mix the errors.

\section{RTL Correspondence}

The gate path has a direct RTL implementation through modules such as \texttt{ewm\_gate\_sbuf\_vec4.sv}. The patch embedding is a logical front-end projection layer and is not the same as the Mamba block internal gate or scan module.

\section{Discussion}

The results show that C++ reference sweeps are useful for identifying quantization-sensitive model components, while \texttt{cpp\_like} versus \texttt{hw\_like} stage comparisons are useful for locating hardware-oriented fixed-point errors. The selective scan result demonstrates that scale design is not a minor implementation detail but a key factor in preserving model behavior.

